\scanpage{45}

of the homomorphism of groupoids $\Pi_{D^*}\longrightarrow\Pi_e$ into a homomorphism which induces an isomorphism on $\pi_0$ and an epimorphism on the $\pi_1$, followed by a homomorphism which is (an epimorphism on $\pi_0$, and for good reason!, and which is) \uline{faithful}. Then the $1$-morphisms from
\[
\Pi'_{D^*}\xrightarrow{\rho'}\Pi'_U\xrightarrow{\varphi'}\Pi_e
\]
\deleted{to}
to
\[
\Pi_{D^*}\xrightarrow{\rho}\Pi_U\xrightarrow{\varphi}\Pi_e
\]
are the \deleted{\unsure{pairs}}\inserted{quintuples} of $2$ functors \inserted{and $3$ isomorphisms of functors}
$(f_{D^*},f_U,\alpha_{D^*U},\alpha_{Ue},K$
giving a diagram with commutation data
\diagram{\begin{tikzcd}[ampersand replacement=\&,column sep=1.05em,row sep=0.75em]
\Pi'_{D^*}\arrow[rrr,"\rho'"]\arrow[ddd,"f_{D^*}"'] \& \& \& \Pi'_U\arrow[rrr,"\varphi'"]\arrow[ddd,"f_U"'] \& \& \& \Pi_e\arrow[ddd,equal] \\
\& \& {} \& \& \& {} \& \\
\& {}\arrow[ur,bend left=25,"\alpha_{D^*U}"] \& \& \& {}\arrow[ur,bend left=25,"\alpha_{Ue}"] \& \& \\
\Pi_{D^*}\arrow[rrr,"\rho"'] \& \& \& \Pi_U\arrow[rrr,"\varphi"'] \& \& \& \Pi_e
\end{tikzcd}},

\margin{moreover, \deleted{\illegible{1}} $\pi_1(\varphi\rho)=\pi_1(\Pi_{D^*}/\Pi_e)$ can be defined as a local system \inserted{of groups (\unsure{commutative|abelian})} on $\Pi_{D^*}$, and we can then \unsure{define|write}
$\boxed{\unsure{\bar k|k}:(\varphi\rho)^*\unsure{T_K|\mathcal T_K}\simeq\pi_1(\varphi\rho)}$.}

\margin{with \uline{one} condition on the induced morphism
$\pi_1(\varphi'\rho')\unsure{\simeq|\longrightarrow}\unsure{f_{D^*}^{\,*}}(\pi_1(\varphi\rho))$,
which modulo the Kummer isomorphisms must be multiplication by a $d$ \unsure{ramification index} (which is a map $\pi_0(\Pi'_{D^*})\longrightarrow\mathbb N^*$). [N.B. multiplication by an \uline{\unsure{integer}} that is \unsure{positive} \illegible{5} constant\ldots]}

if we have two such $1$-morphisms
\[
f=(f_{D^*},f_U,\alpha_{D^*U},\alpha_{Ue})
\quad\text{and}\quad
g=(g_{D^*},g_U,\beta_{D^*U},\beta_{Ue})
\]
a \unsure{$0$-morphism|homomorphism} from $f$ to $g$ is formed of a pair of $\mu_{D^*},\mu_U$ of isomorphisms of functors
\[
\mu_{D^*}:f_{D^*}\simeq g_{D^*},\qquad
\mu_U:f_U\longrightarrow g_U
\]
compatible with $\alpha_{D^*U},\inserted{\beta_{D^*U},}$ and with $\alpha_{Ue},\inserted{\beta_{Ue},}$ \deleted{and with}
\margin{44}

\scanpage{46}
\mspage{21}

\deleted{$k,k'$ modulo multiplication by a ramification index} (two compatibility conditions on the two squares).

If we \deleted{assume} restrict ourselves to connected $\Pi_U,\Pi'_U$, we find by choosing as above an object \inserted{$\widetilde U$} of $\Pi_U$ and an isomorphism
\[
\boxed{\varphi(\widetilde U)\overset{\lambda}{\simeq}\Omega}
\]
(where $\Omega=\operatorname{Spec}(\bar K)$ is the reference object of $\Pi_e\ldots$), a group $\mathcal E$ and a homomorphism
\[
\mathcal E\xrightarrow{\varphi}\Gamma
\]
(which is changed by an inner automorphism if we change $\lambda$, and $\mathcal E$ itself replaced by an isomorphic group, \unsure{with isomorphism} defined \unsure{modulo} inner isomorphisms, if we change the reference object $\widetilde U$). Likewise, choosing a $\widetilde U_i$ in each component $i\in\pi_0(\Pi_{D^*})$, and isomorphisms
\[
\lambda_i:\rho(\widetilde U_i)\simeq\widetilde U,
\]
we find groups $\mathcal E_i$ and homomorphisms
\[
\mathcal E_i\xrightarrow{\rho_i}\mathcal E.
\]
If we change $\lambda_i$, $\rho_i$ is changed by an inner automorphism of $\mathcal E$. If we change $\widetilde U_i$, $\mathcal E_i$ is replaced by an isomorphic group, \inserted{the isomorphism} defined up to an inner automorphism.
\margin{45}

\scanpage{47}

Let us consider the composites
\diagram{\begin{tikzcd}[ampersand replacement=\&]
\mathcal E_i\arrow[r]\arrow[rr,bend left=30,"p_i"]\&\mathcal E\arrow[r,"p"']\&\Gamma
\end{tikzcd}}
whence kernel, \deleted{\illegible{1}} $\unsure{K|k}$ is defined by a system of isomorphisms
\[
T(\bar K)\overset{k_i}{\simeq}\operatorname{Ker}p_i\qquad(i\in I)
\]
fitting into a family of exact sequences
\[
1\longrightarrow T(\bar K)\xrightarrow{k_i}\mathcal E_i\xrightarrow{p_i}\Gamma
\]
(N.B. the image of $p_i$ is an open subgroup)
$k_i$ being compatible with the actions of $\mathcal E_i$, when we let $\mathcal E_i$ act on $T(\bar K)$ through the action of $\Gamma$ on $T(\bar K)$, and on itself by inner automorphisms. Introducing also $\pi=\operatorname{Ker}p$, we thus find homomorphisms of exact sequences
\diagram{$(D)$\quad\begin{tikzcd}[ampersand replacement=\&]
1\arrow[r]\&T(\bar K)\arrow[r,"k_i"]\arrow[d,"\varepsilon_i\,\deleted{\illegible{1}}"']\&\mathcal E_i\arrow[r,"p_i"]\arrow[d,"\rho_i"]\&\Gamma\arrow[d,equal]\\
1\arrow[r]\&\pi\arrow[r,"\unsure{k|\kappa}"']\&\mathcal E\arrow[r,"p"']\&\Gamma
\end{tikzcd}\qquad$i\in I$}

\deleted{Let us consider}
We can say that $\Pi_U$ defines an ``outer group'' $[\mathcal E]$, and $\Pi_U\longrightarrow\Pi_e$ a homomorphism of outer groups
\[
[\mathcal E]\longrightarrow[\Gamma],
\]
likewise $\Pi_{D^*}$ defines a system of ``outer groups'' $[\mathcal E_i]$, and \deleted{\unsure{of the homomorphism}} $\Pi_{D^*}\longrightarrow\Pi_U$
\margin{46}

\scanpage{48}
\mspage{22}

a system of outer homomorphisms
\[
[\mathcal E_i]\longrightarrow[\mathcal E],
\]
\deleted{finally} but how in terms of outer groups to express the Kummer data $k_i$?

Returning to the \inserted{system of} diagram $D$ (relative to an index set $I$.) and to an analogous $D'$ (with index set $I'$)
\diagram{$D'$\quad\begin{tikzcd}[ampersand replacement=\&]
1\arrow[r]\&T(\bar K)\arrow[r,"k'_{i'}"]\arrow[d,"\varepsilon'_{i'}"']\&\mathcal E'_{i'}\arrow[r,"p'_{i'}"]\arrow[d,"\rho'_{i'}"]\&\Gamma\arrow[d,equal]\\
1\arrow[r]\&\unsure{\pi|\pi'}\arrow[r,"\unsure{k'|\kappa'}"']\&\mathcal E'\arrow[r,"p'"']\&\Gamma
\end{tikzcd}\qquad$i'\in\unsure{I|I'}$}
a $1$-morphism \inserted{$f$} from $D'$ to $D$ is made explicit by

$a)$ A homomorphism
\[
\boxed{\boxed{f_{\mathcal E}:\mathcal E'\longrightarrow\mathcal E}}
\]
\margin{describes the \deleted{\unsure{corresponding}} functor $f_U$}
[\deleted{\unsure{depends}}\inserted{is made explicit} in terms of the choice of an isomorphism
\[
\nu:f(\widetilde U')\simeq\widetilde U
\]
[another choice modifies $f_{\mathcal E}$ by an $\operatorname{int}(\beta)$, $\beta\in\mathcal E$)]

$b)$ A \deleted{homomorphism} map $\tau:I'\longrightarrow I$, and for every $i'\in I'$, putting $i=\tau(i')$, a \deleted{map}\inserted{group homomorphism}
\[
\boxed{\mathcal E'_{i'}\xrightarrow{f_{i'}}\mathcal E_i}
\]
\margin{describes the functor $f_{D^*}$}
\margin{47}

\scanpage{49}

[is made explicit in terms of the choice of an isomorphism
\[
f_{D^*}(\widetilde U'_{i'})\overset{\nu_{i'}}{\simeq}\widetilde U_i,
\]
and modified by $\operatorname{int}(\beta_{i'})$, $\beta_{i'}\in\mathcal E_i$, if we change $\nu_{i'}$]

$c)$ A commutation datum for
\diagram{\begin{tikzcd}[ampersand replacement=\&,column sep=0.9em,row sep=0.7em]
\mathcal E'\arrow[rrr,"p'"]\arrow[ddd,"f_{\mathcal E}"'] \& \& \& \unsure{\Gamma|\Gamma'}\arrow[ddd,equal] \\
\& \& {} \& \\
\& {}\arrow[ur,bend left=25,"{\deleted{\unsure{\beta}}\alpha}"] \& \& \\
\mathcal E\arrow[rrr,"p"'] \& \& \& \Gamma
\end{tikzcd}\quad$\boxed{\boxed{\deleted{\unsure{\beta}}\alpha\in\Gamma}}$}
\margin{$\alpha$ describes $\alpha_{Ue}$}
i.e.
\[
\deleted{f_{\mathcal E}p'=\operatorname{int}(\unsure{\beta})p}
\qquad\boxed{pf_{\mathcal E}=\operatorname{int}(\alpha)p'}
\]

$d)$ $\forall i'[\in\unsure{I|I'}]$ a commutation datum for
\diagram{\begin{tikzcd}[ampersand replacement=\&,column sep=0.9em,row sep=0.7em]
\mathcal E'_{i'}\arrow[rrr,"\rho'_{i'}"]\arrow[ddd,"f_{i'}"'] \& \& \& \mathcal E'\arrow[ddd,"f_{\mathcal E}"] \\
\& \& {} \& \\
\& {}\arrow[ur,bend left=25,"\alpha_{i'}"] \& \& \\
\mathcal E_i\arrow[rrr,"\rho_i"'] \& \& \& \mathcal E
\end{tikzcd}\quad$\boxed{\alpha_{i'}\in\mathcal E}\quad(i'\in I')$}
\margin{$(\alpha_{i'})$ describes $\alpha_{D^*U}$}
i.e.:
\[
\boxed{\rho_i f_{i'}=\operatorname{int}(\alpha_{i'})f_{\mathcal E}\rho'_{i'}}
\]

Let us note that $c,d$ together define a commutation datum
\diagram{\begin{tikzcd}[ampersand replacement=\&,column sep=0.9em,row sep=0.7em]
\mathcal E'_{i'}\arrow[rrr,"p'_{i'}"]\arrow[ddd,"f_{i'}"'] \& \& \& \Gamma\arrow[ddd,equal,"{\deleted{\operatorname{int}(\unsure{\gamma})}}"] \\
\& \& {} \& \\
\& {}\arrow[ur,bend left=25,"{\operatorname{int}(\beta_{i'})}"] \& \& \\
\mathcal E_i\arrow[rrr,"p_i"'] \& \& \& \Gamma
\end{tikzcd}}
with $\beta_{i'}=p(\alpha_{i'})\alpha$
i.e.
\[
p_i f_{i'}=\operatorname{int}(\beta_{i'})p'_{i'}
\]
\margin{48}

\scanpage{50}
\mspage{23}

i.e. we have commutativity in
\diagram{\begin{tikzcd}[ampersand replacement=\&]
\mathcal E'_{i'}\arrow[r,"p'_{i'}"]\arrow[d,"f_{i'}"']\&\Gamma\arrow[d,"\operatorname{int}(\beta_{i'})"]\\
\mathcal E_i\arrow[r,"p_i"']\&\Gamma
\end{tikzcd}}
so $f_{i'}$ induces a homomorphism of exact sequences
\diagram{\begin{tikzcd}[ampersand replacement=\&]
1\arrow[r]\&T(\bar K)\arrow[r,"k'_{i'}"]\arrow[d,"f^0_{i'}"']\&\mathcal E'_{i'}\arrow[r,"\unsure{p'_{i'}|f'_{i'}}"]\arrow[d,"f_{i'}"]\&\Gamma\arrow[d,"\operatorname{int}(\beta_{i'})"]\\
1\arrow[r]\&T(\bar K)\arrow[r,"k_i"']\&\mathcal E_i\arrow[r,"p_i"']\&\Gamma
\end{tikzcd}}
\margin{795}
and we must express the compatibility of $f^0_{i'}$ with the ramification index $d_{i'}\in\mathbb N$, we find
\[
\boxed{f^0_{i'}=d_{i'}\chi(\beta_{i'})\operatorname{id}_{T(\bar K)}}
\]
where
\[
\chi:\Gamma\longrightarrow\widehat{\mathbb Z}^{\,*}
\]
is the canonical character.

This being so, let us determine, for a $1$-morphism
\[
f=(f_{\mathcal E},\inserted{\tau,}(f_{i'}),\alpha,(\unsure{\alpha_{i'}|\alpha_i}))
\]
and another
\[
f'=(f'_{\mathcal E},\tau',(f'_{i'}),\alpha',(\unsure{\alpha'_{i'}|\alpha'_i})),
\]
the $0$-morphisms from one to the other.
\margin{49}